%==============================================================================
\section{Discussion}
\label{sec:discussion}

\subsection{What typed verdicts buy}
\label{sec:disc-verdicts}

An inference engine asked for the Gibbs energy of 2{,}914 cytosolic reactions returns 75 values and silence.
The kernel returns 2{,}914 verdicts, and the silence is resolved by cause.
For 2{,}244 reactions the question is outside what TECRDB can state; for 398 no measurement exists; for 197 the reaction cannot even be put to TECRDB, because a participant has no identifier it uses.
These three populations need different remedies.
The first needs a different kind of model, such as the carrier/slot decomposition, which recovered 96.
The second needs measurement or estimation.
The third needs cross-references, and I1 (\cref{sec:res-db}) shows how far those go: MetaNetX moved 33 reactions from the third population into the second and none into the answered set.
The gate letters made that movement visible as a change from \textsf{open} to $-$.

The payload rule (V1) did real work.
A missing $\DG$ never became zero, a missing $K_m$ never became a typical value, and an undecided direction never became the curated one.
Every downstream count in this paper, from 156 concentration-blocked directions to 3 answered branch points, inherits that discipline.
The rule also makes failures cheap to diagnose: the S1 parser fault surfaced as an implausibly low fraction rather than as silently wrong slot energies, because the slot check reads curated features directly.

Basis tags and ground verdicts answer a different question: not ``is there a value?'' but ``what kind of value is it?''.
The AlphaFold contact graphs are icons, models of their object; A1 shows that on 60 proteins they agree with indices, crystal structures, at a median Jaccard index of 0.90 (\cref{fig:structures}).
That agreement licenses using them where no crystal structure exists, but it does not turn them into measurements, and the ground verdict keeps saying so.
The falsifiability guard (\cref{prop:falsifiable}) is the same idea applied to cycles: a consistency check that cannot fail is refused, not passed.

\subsection{Concentrations are the binding constraint}
\label{sec:disc-conc}

Three independent results point to the same missing datum.
M2b found that for 156 of 199 undecided reactions the medium fixes the sign of $\DGm$, so only concentrations remain.
R1 found that the curated directions which disagree with the 1\,mM sign are those of reactions running uphill at 1\,mM, such as triose-phosphate isomerase and glyceraldehyde-3-phosphate dehydrogenase, which glycolysis pulls forward by removing product.
The completion survey found that fixing those directions adds 11\% more unique completions (\cref{sec:res-complete}).
Absolute concentrations measured in mammalian cells \citep{park2016concentrations} are therefore the single most valuable addition to the kernel.
The loader for them exists (\cref{sec:loaders}); its file does not yet.

The same result reframes what a ``direction'' annotation is.
A curated direction is a statement about the physiological concentration regime, not about the reaction.
The kernel's direction verdict makes the dependence explicit: its blocker names the concentration window, and its remedy names the measurement.

\subsection{Enzymes are not dilute}
\label{sec:disc-locality}

The chromatographic reading of the cytosol gives two numbers per contact, and they point in opposite directions.
From the metabolite's side, retention is negligible: the median $k$ is 0.001--0.003, so binding barely changes a metabolite's capacitance.
From the enzyme's side, contacts are non-local (\cref{sec:res-kinetics}): the median enzyme departs from its dilute limit at 2.4\,\textmu M of substrate, and 221 of 223 depart by more than 10\% at 1\,mM.
The consequence for modelling is concrete.
Enzyme--substrate states cannot be folded into metabolite capacitances; they must stay explicit, with conserved enzyme totals, which is what the two-rung expansion provides.
The failure of the glucose divider in E3 is the same fact seen from the flux side: shares computed in the dilute limit are wrong where the enzyme is saturated.

\subsection{What failed, and why that matters}
\label{sec:disc-negative}

The coupled-oscillator account of folding failed every test it could fail.
Its sequence coupling carries less contact information than a hydrophobicity product (F1).
Its eigenvalue count misses even with the true contact graph as coupling (F2).
Its phase-lock groups lose to contiguous blocks (F3), and its dynamics add nothing to the static coupling they run on ($p=0.14$).
F4 appeared to rescue a weak effect by attributing it to disorder, and D1 then removed the rescue, because the effect persists outside measured disordered regions.
The pattern argues against adopting the account as a structure source and for keeping AlphaFold's counted contact graph, whose agreement with experiment A1 established (\cref{fig:structures}).

M3 failed in an instructive way.
The two reactions whose direction the kernel decides clear zero by 0.28\kJ; a crowding shift of 0.92\kJ\ removes both.
A direction verdict that a plausible change of convention can overturn is not robust, and the kernel reports it as such.

G1 held, but only 23 of its 324 cases could have failed (\cref{sec:res-db}), and the paper says so.
A check is only as strong as its discriminating cases, which is why the check markers distinguish \emph{pass} from \emph{non-discriminating}.

\subsection{Limitations and threats to validity}
\label{sec:limitations}

\paragraph{Scope of the corpus.}
The corpus is defined by Reactome's compartment annotations.
L1 found 78 reactions none of whose catalysts has cytosol support, including secreted fibrinogen chains; these are likely annotation errors that the filter inherited.
Reactions that span compartments, including all transport, are excluded by construction.

\paragraph{Coverage of measured thermodynamics.}
Only 61 distinct TECRDB reactions match the corpus, and they close no cycle, so the cycle test is \vd{certified-empty} rather than informative.
The wider TECRDB set is inconsistent at $\chi^2/\nu=4.11$, so its quoted uncertainties are optimistic by a factor of about two.
The pooled $\sigma$ of 2.68\kJ\ is applied uniformly and does not capture reaction-specific error.

\paragraph{Conventions are supplied, not measured.}
The medium window, the concentration window, the ambient list, the conversion from parts per million to molar and the accessible-volume range are typical cytosolic values, not measurements for any cell type.
Each is named on every value it affects, and each enters the stage digest, but a different choice changes the numbers.
Abundances are whole-organism integrated values except in the HeLa comparison.

\paragraph{Temperature.}
Transforms are evaluated at 298.15\,K because TECRDB values are not temperature-corrected.
The concentration term of the direction verdict and the junction law use $RT$ at 310\,K (2.577\kJ).
The mismatch scales $RT\ln Q$ by about 4\%, and it is not negligible.
Evaluating the concentration term at 298.15\,K as well would decide eight further reactions forward.
All eight are cAMP hydrolysis by phosphodiesterases, with Hess-derived values, and they clear zero by 0.8\kJ.
Like the two pyrophosphatase verdicts, these sit within a convention's reach of the boundary.
A consistent treatment would correct the standard values to 310\,K before any of them is relied on.

\paragraph{Multiple testing.}
The 23 claims were tested at $\alpha=0.05$ without correction, as declared.
Most decisive outcomes are far from the threshold, for example C1 ($p=2.7\times10^{-18}$) and the shuffle tests ($p=0.005$, the smallest attainable with 200 permutations).
F4a ($p=0.012$) and the pLDDT difference within it ($p=0.043$) would not survive a Bonferroni correction over 23 claims, and should be read as suggestive.

\paragraph{Single runs and implementation faults.}
Each stage was run once on fixed data.
One implementation fault was found during the work and is reported with its effect: the UniProt parser initially truncated notes at semicolons, which lowered S1 from 81.6\% to 19\% until it was fixed.
The test suite (\cref{app:tests}) guards each such fault found so far, but it cannot exclude faults not yet found.

\paragraph{Missing data sources.}
Two sources the design anticipates could not be obtained under their terms.
The supplementary concentration tables of \citet{park2016concentrations} require a publisher login, and the Human Metabolome Database refused scripted downloads \citep{wishart2022hmdb}.
BRENDA requires a licence to be accepted interactively \citep{chang2021brenda}, and PhosphoSitePlus requires registration under a non-commercial licence \citep{hornbeck2015psp}.
Their loaders are in place, and their verdicts are \vd{absent} until the files are supplied.

%==============================================================================
\section{Conclusion}
\label{sec:conclusion}

A knowledge graph of biochemistry need not infer.
When its edges are element laws and its nodes are solutions, its structure yields rank, conservation laws, blocked reactions and cycle constraints without any parameter.
When its answers are typed verdicts with non-launderable provenance, every gap is named by cause, and every remedy is explicit.
Applied to the human cytosol, the kernel shows how little of the network public measurements determine (\cref{tab:outcomes,fig:coverage}): 201 of 2{,}914 reaction energies, 2 robust directions, 3 flux dividers.
It also shows precisely what would determine more (\cref{sec:disc-conc}).
Concentrations would decide most directions; measurement, not identifiers, limits thermodynamic coverage; enzyme saturation, not metabolite sequestration, governs the contacts.
The pre-registered audit kept the kernel honest about all of this, including the claims that failed.
